\documentclass[10pt]{article}
\usepackage[letterpaper,margin=0.7in]{geometry}
\usepackage[T1]{fontenc}
\usepackage{helvet}
\renewcommand{\familydefault}{\sfdefault}
\usepackage{enumitem}
\usepackage{tabularx}
\usepackage{booktabs}
\usepackage{xcolor}
\usepackage[hidelinks]{hyperref}
\definecolor{accent}{RGB}{19,28,45}
\setlength{\parindent}{0pt}
\setlength{\parskip}{4pt}
\renewcommand{\arraystretch}{1.25}
\pagestyle{empty}
\newcommand{\sect}[1]{\vspace{5pt}{\large\bfseries\color{accent}#1}\vspace{2pt}\hrule\vspace{3pt}}
\newcolumntype{Y}{>{\raggedright\arraybackslash}X}

\begin{document}

{\LARGE\bfseries Go-to-market executive summary, in five layers}\\[3pt]
{\color{accent}\textbf{aimez}} \quad strategy summary \quad June 2026

Each layer stands alone in conversation. Start at whichever layer the listener cares about; the layers below are the support.

\sect{Layer 1: Product}
A measured decision surface for pedestrian movement in Manhattan. A street graph where every block carries validated layers (street pressure, tree canopy, surface heat) and an engine that returns five comparable routes with per-block statistics and confidence in one call. Consumers see routes; the asset is the surface underneath, which knows where people walk and how each block feels. Proof: the Grand Central to Carnegie Hall case (\texttt{fig\_case\_gc\_carnegie}).

\sect{Layer 2: Moat}
The validation and provenance discipline, not the routing code. Every layer carries a manifest with input hashes; every claim has an allowed and not-allowed list; audits are cheap enough to run the same morning a doubt is raised. Live proof this week: a tree-level audit caught the canopy layer crediting side-street trees to avenue corners, withdrew a published claim and replaced the mechanism within one day (\texttt{fig\_case\_upwind}). A competitor can re-weight public data; they cannot shortcut the discipline that catches that class of error.

\sect{Layer 3: Beachhead, out-of-home advertising}
OOH sells foot-traffic exposure as its entire currency and buys data products, not apps. The offer: a comfort-weighted exposure index per panel location with a confidence score and a provenance trail. Incumbent counts say how many; the surface explains why a block draws or repels walkers and how long they linger. Pilot: one corridor, one metric, one media owner, success defined as the index reordering inventory in a way the current data missed. Mechanics proof: the Village dog-park case (\texttt{fig\_case\_village}), flip threshold about 11 DBH-inches per 100 m.

\sect{Layer 4: Adjacent markets}
In priority order behind the beachhead: outdoor-workforce heat compliance under tightening occupational rules, municipal resilience and transportation offices that buy static heat reports today, real estate walkability indices, last-mile logistics routing. Each reuses the same surface with a different cost weighting; none needs a consumer user base.

\sect{Layer 5: Motion and asks}
Motion: paid pilot study, then a per-market annual data license. Land one media owner, expand to brand-side buyers. Research leg: NSF SBIR Phase I pitch submitted. Asks in order: go-to-market judgment, adjacent paid work near decision systems, a milestone-gated investor look after the pitch lands and a pilot signs.

\sect{Layer summary}
\begin{tabularx}{\textwidth}{@{}lYl@{}}
\toprule
\textbf{Layer} & \textbf{The one line} & \textbf{Proof point} \\
\midrule
Product & A measured decision surface for foot traffic & GC to Carnegie case \\
Moat & Validation and provenance, not the routing code & The canopy audit catch \\
Beachhead & A comfort-weighted exposure index for OOH & The Village dog-park case \\
Adjacent & Same surface, different cost weighting & Compliance, cities, real estate \\
Motion & Paid pilot then license; three layered asks & NSF pitch submitted \\
\bottomrule
\end{tabularx}

\sect{Honest caveats}
OOH measurement is standards-driven, so an independent index has to be benchmarked before owners adopt it. Comfort-to-dwell is a hypothesis the first pilot exists to test. Mobile-location panels out-scale us, so the pitch stays on explanatory quality, not raw coverage. All heat readings are surface temperature until street-level measurement exists.

\end{document}
